\documentclass[a4paper,12pt]{article}
\usepackage[T2A]{fontenc}
\usepackage[utf8]{inputenc}
\usepackage[russian]{babel}
\usepackage{amsmath,amssymb,mathtools}
\usepackage{icomma}
\usepackage[left=22mm,right=22mm,top=21mm,bottom=22mm]{geometry}
\usepackage{microtype}
\usepackage{setspace}
\usepackage{parskip}
\usepackage{enumitem}
\usepackage{booktabs,tabularx,array}
\usepackage{xcolor}
\usepackage{tikz}
\usetikzlibrary{arrows.meta,positioning,shapes.geometric,calc}
\usepackage{fancyhdr}
\usepackage{hyperref}
\definecolor{accent}{HTML}{006F73}
\definecolor{darkaccent}{HTML}{004D51}
\definecolor{textgray}{HTML}{333333}
\definecolor{muted}{HTML}{717777}
\hypersetup{colorlinks=true,linkcolor=darkaccent,urlcolor=accent,pdfauthor={Артём Александрович Лёвин},pdftitle={Чек-лист. Область определения, иррациональные выражения и множество значений}}
\onehalfspacing
\setlength{\parindent}{0pt}
\setlength{\parskip}{5pt}
\setlist[enumerate]{leftmargin=7mm,itemsep=7pt,topsep=6pt}
\pagestyle{fancy}
\fancyhf{}
\fancyhead[L]{\small\color{muted}ОДЗ и множество значений функций}
\fancyhead[R]{\small\color{muted}Анна Трапезникова}
\fancyfoot[C]{\small\color{muted}\thepage}
\renewcommand{\headrulewidth}{0.4pt}
\renewcommand{\headrule}{\hbox to\headwidth{\color{accent}\leaders\hrule height \headrulewidth\hfill}}
\setlength{\headheight}{25pt}
\newcommand{\sectionhead}[2]{\vspace{7pt}{\Large\color{darkaccent}#1. #2\par}\vspace{3pt}{\color{accent}\rule{\linewidth}{0.45pt}}\vspace{6pt}}
\newcommand{\subhead}[1]{\vspace{5pt}{\color{darkaccent}\large #1\par}\vspace{2pt}}
\newcommand{\idea}[1]{\textcolor{darkaccent}{\textbf{Главная идея.}} #1}
\newcommand{\R}{\mathbb{R}}
\begin{document}
\thispagestyle{empty}
\begin{center}
\vspace*{2.6cm}
{\large\color{darkaccent}\MakeUppercase{Математика}}\par
\vspace{1.5cm}
{\LARGE\color{textgray}Чек-лист}\par
\vspace{0.9cm}
{\LARGE\color{accent}Область допустимых значений\\[2mm] и множество значений функций}\par
\vspace{1.15cm}
{\large\color{textgray}Алгебра · 10 класс}\par
\vfill
{\large Анна Трапезникова}\par
\vspace{0.65cm}
{\large 3 ноября 2026 г.}\par
\vspace{1.0cm}
{\normalsize Лёвин Артём Александрович\\[-1mm]
эксклюзивно для Анны Трапезниковой}\par
\vspace{1.1cm}
\begin{tikzpicture}
\draw[accent,line width=1pt] (0,0) .. controls (1.6,.7) and (3.15,-.6) .. (4.6,0)
 .. controls (6.2,.65) and (7.7,-.65) .. (9.4,0);
\fill[accent] (4.6,0) circle (1.6pt);
\end{tikzpicture}
\end{center}
\clearpage

\sectionhead{1}{Зачем нужна область допустимых значений}
\idea{Перед преобразованиями определите, при каких $x$ исходное выражение имеет смысл.}

Для многочлена, например $x^2-4$, допустимы \emph{все} действительные числа: $D=\R$. Ограничения возникают, когда в выражении есть деление, корень чётной степени и некоторые другие операции.

\sectionhead{2}{Три правила, которые нужно помнить}
\renewcommand{\arraystretch}{1.45}
\begin{tabularx}{\linewidth}{@{}p{0.30\linewidth}X@{}}
\toprule
\textcolor{darkaccent}{\textbf{Что встретилось}} & \textcolor{darkaccent}{\textbf{Что записать}} \\
\midrule
$\sqrt{A(x)}$ & $A(x)\geqslant0$ — подкоренное выражение неотрицательно. \\
$\dfrac{B(x)}{C(x)}$ & $C(x)\ne0$ — знаменатель не равен нулю. \\
Несколько условий & Все ограничения должны выполняться \emph{одновременно}: находите их пересечение. \\
\bottomrule
\end{tabularx}

\subhead{Пример 1. Корень из квадратного трёхчлена}
\[
\sqrt{x^2-4}:\qquad x^2-4\geqslant0\quad\Longleftrightarrow\quad
x\leqslant-2\ \text{или}\ x\geqslant2.
\]
Значит, $D=(-\infty;-2]\cup[2;+\infty)$. Точки $-2$ и $2$ включены: $\sqrt0=0$.

\subhead{Пример 2. Корень с обратным ограничением}
\[
\sqrt{9-x^2}:\qquad 9-x^2\geqslant0
\quad\Longleftrightarrow\quad -3\leqslant x\leqslant3.
\]
Получаем $D=[-3;3]$. Обратите внимание: направление неравенства влияет на итоговые интервалы.

\subhead{Пример 3. Знаменатель}
\[
\frac{1}{x^2-5}:\qquad x^2-5\ne0
\quad\Longleftrightarrow\quad x\ne\pm\sqrt5.
\]
Здесь допустимы все действительные числа, кроме двух указанных точек.

\clearpage
\sectionhead{3}{Как объединять ограничения}
Рассмотрим выражение
\[
F(x)=\frac{\sqrt{x^2-4}}{x^2-5}.
\]
Одновременно нужны два условия:
\[
\begin{cases}
x^2-4\geqslant0,\\
x^2-5\ne0.
\end{cases}
\qquad\Longleftrightarrow\qquad
\begin{cases}
x\leqslant-2\ \text{или}\ x\geqslant2,\\
x\ne-\sqrt5,\quad x\ne\sqrt5.
\end{cases}
\]
Поэтому
\[
\boxed{D=\bigl((-\infty;-2]\cup[2;+\infty)\bigr)\setminus\{-\sqrt5;\sqrt5\}.}
\]
Знаки $\geqslant$ и $\ne$ отвечают за \emph{разные причины} ограничения. Нельзя заменить их одним общим условием «всё больше нуля».

\sectionhead{4}{Уравнение с корнем: изолируйте радикал}
Чтобы решить уравнение $\sqrt{A(x)}=B(x)$, пользуйтесь равносильной системой:
\[
\boxed{\begin{cases}
A(x)\geqslant0,\\ B(x)\geqslant0,\\ A(x)=B^2(x).
\end{cases}}
\]
Условие $B(x)\geqslant0$ принципиально: квадратный корень не бывает отрицательным. После возведения в квадрат следует проверить все найденные значения в исходном уравнении.

\subhead{Пример 4. Замена $t=x^2$}
Решим $\sqrt{x^2-4}=6-x^2$.

Сначала вводим $t=x^2$. Тогда $t\geqslant0$, а из условий на обе части уравнения получаем $4\leqslant t\leqslant6$.
\[
\sqrt{t-4}=6-t
\quad\Longrightarrow\quad t-4=(6-t)^2
\quad\Longrightarrow\quad t^2-13t+40=0.
\]
\[
(t-5)(t-8)=0\quad\Longrightarrow\quad t=5\ \text{или}\ t=8.
\]
Число $8$ не подходит: нарушает $t\leqslant6$. При $t=5$ исходное равенство верно: $\sqrt1=1$.
\[
\boxed{x=-\sqrt5\quad\text{или}\quad x=\sqrt5.}
\]

\clearpage
\sectionhead{5}{Как найти нули сложного знаменателя}
Разберём выражение, объединяющее идеи занятия:
\[
G(x)=\frac{\bigl(\sqrt{x^2-4}-1\bigr)\bigl(\sqrt{x^2-4}+5\bigr)}{x^2-6+\sqrt{x^2-4}}.
\]
Корень требует $x^2-4\geqslant0$. Также нужно исключить решения уравнения
\[
x^2-6+\sqrt{x^2-4}=0
\quad\Longleftrightarrow\quad
\sqrt{x^2-4}=6-x^2.
\]
Оно уже решено в примере 4: знаменатель равен нулю \emph{только} при $x=\pm\sqrt5$. Числитель дополнительных ограничений не создаёт. Поэтому
\[
\boxed{D(G)=\bigl((-\infty;-2]\cup[2;+\infty)\bigr)\setminus\{-\sqrt5;\sqrt5\}.}
\]
\subhead{Почему важно проверять корни}
Уравнение $A(x)=B(x)$ после возведения в квадрат может дать лишние значения. Например, для $t=8$ получаем $\sqrt{8-4}=2$, а $6-8=-2$. Квадраты обеих частей совпадают, исходные части — нет.

\sectionhead{6}{Неравенство с модулем}
При $a\geqslant0$ действует правило
\[
\boxed{|H(x)|\leqslant a\quad\Longleftrightarrow\quad -a\leqslant H(x)\leqslant a.}
\]
\subhead{Пример 5}
\[
\begin{aligned}
|x^2-5|\leqslant4
&\quad\Longleftrightarrow\quad -4\leqslant x^2-5\leqslant4\\
&\quad\Longleftrightarrow\quad 1\leqslant x^2\leqslant9.
\end{aligned}
\]
Это означает $1\leqslant |x|\leqslant3$, следовательно,
\[
\boxed{x\in[-3;-1]\cup[1;3].}
\]


\clearpage
\sectionhead{7}{Область определения и множество значений}
В предыдущих разделах мы искали \textbf{область определения}~$D(f)$: все значения $x$, при которых выражение $f(x)$ имеет смысл.

Теперь вопрос другой: \textbf{множество значений}~$E(f)$ --- все числа $y$, которые функция действительно может принимать. Для этого изучаем уравнение $f(x)=y$: при каких $y$ существует хотя бы одно допустимое $x$?

Например, у функции $f(x)=x^2+1$ область определения $D(f)=\mathbb R$, а множество значений $E(f)=[1;+\infty)$. Причина: $x^2\geqslant0$ при любом $x$, и значение $1$ достигается при $x=0$.

\subhead{Пример 6. Оценка через неотрицательный квадрат}
Найдём множество значений функции, разобранной на доске:
\[
f(x)=\frac{2x}{x^2+1}.
\]
Знаменатель всегда положителен: $x^2+1\geqslant1>0$. Значит, $D(f)=\mathbb R$.

\textcolor{darkaccent}{\textbf{Шаг 1.}} Используем тот факт, что квадрат любого действительного числа неотрицателен:
\[
(|x|-1)^2\geqslant0.
\]
Раскроем скобки и учтём, что $|x|^2=x^2$:
\[
x^2-2|x|+1\geqslant0
\quad\Longleftrightarrow\quad x^2+1\geqslant2|x|.
\]
\textcolor{darkaccent}{\textbf{Шаг 2.}} Делим на $x^2+1>0$, сохраняя знак неравенства:
\[
\frac{2|x|}{x^2+1}\leqslant1
\quad\Longleftrightarrow\quad |f(x)|\leqslant1.
\]
По определению модуля
\[
|f(x)|\leqslant1
\quad\Longleftrightarrow\quad
\begin{cases}f(x)\leqslant1,\\f(x)\geqslant-1,\end{cases}
\quad\Longleftrightarrow\quad -1\leqslant f(x)\leqslant1.
\]

\textcolor{darkaccent}{\textbf{Шаг 3.}} Проверяем крайние значения:
\[
f(1)=\frac{2}{2}=1,\qquad
f(-1)=\frac{-2}{2}=-1.
\]
Значит, обе границы входят в ответ. Чтобы убедиться, что внутри промежутка нет пропусков, для произвольного $y$ рассмотрим $yx^2-2x+y=0$. При $y\ne0$ его дискриминант равен $4(1-y^2)\geqslant0$ для всех $y\in[-1;1]$; при $y=0$ подходит $x=0$. Поэтому
\[
\boxed{E(f)=[-1;1].}
\]
\idea{Неотрицательный квадрат часто позволяет получить точные границы дробно-рациональной функции без производной.}

\sectionhead{8}{Множество значений через дискриминант}
\subhead{Пример 7. Второе задание с доски}
Найдём множество значений функции
\[
f(x)=\frac{4x+8}{x^2+5}.
\]
Знаменатель $x^2+5>0$ для всех $x$. Положим $y=f(x)$ и выясним, при каких $y$ можно найти действительное $x$:
\[
y=\frac{4x+8}{x^2+5}
\quad\Longleftrightarrow\quad y(x^2+5)=4x+8.
\]
Переносим всё в одну сторону и группируем по степеням $x$:
\[
\boxed{yx^2-4x+5y-8=0.}
\]
\textcolor{darkaccent}{\textbf{Шаг 1.}} Временно считаем $y\ne0$. Это квадратное уравнение относительно $x$ с коэффициентами
\[
a=y,\qquad b=-4,\qquad c=5y-8.
\]
Оно имеет действительные корни тогда и только тогда, когда дискриминант неотрицателен:
\[
\begin{aligned}
D_x&=(-4)^2-4y(5y-8)\\
&=16-20y^2+32y\geqslant0.
\end{aligned}
\]
\textcolor{darkaccent}{\textbf{Шаг 2.}} Преобразуем неравенство (помним о смене знака при умножении на отрицательное число):
\[
16-20y^2+32y\geqslant0
\quad\Longleftrightarrow\quad 5y^2-8y-4\leqslant0.
\]
Корни трёхчлена:
\[
y_{1,2}=\frac{8\pm\sqrt{64+80}}{10}
=\frac{8\pm12}{10},\qquad y_1=-\frac25,\quad y_2=2.
\]
Парабола $5y^2-8y-4$ направлена ветвями вверх, поэтому выражение не положительно \emph{между корнями, включая их}:
\[
-\frac25\leqslant y\leqslant2.
\]
\textcolor{darkaccent}{\textbf{Шаг 3.}} При $y=0$ исходное уравнение становится линейным: $-4x-8=0$, откуда $x=-2$. Поэтому нулевое значение тоже достижимо.

Границы действительно достигаются: $f(1)=2$, а $f(-5)=-\frac25$. Следовательно,
\[
\boxed{E(f)=\left[-\frac25;2\right].}
\]

\subhead{Как действовать в аналогичных задачах}
\begin{enumerate}[leftmargin=7mm,itemsep=2pt,topsep=2pt]
\item Определите $D(f)$ и исключите запрещённые значения $x$.
\item Запишите $y=f(x)$ и получите уравнение относительно $x$.
\item Для квадратного уравнения потребуйте $D_x\geqslant0$ и решите неравенство для $y$.
\item Если старший коэффициент зависит от $y$, отдельно проверьте значения, при которых он обращается в нуль.
\item Проверьте достижимость границ и запишите $E(f)$.
\end{enumerate}
\textbf{Важно.} Неравенство $D_x\geqslant0$ отвечает на вопрос о существовании действительного корня $x$. После умножения на знаменатель важно помнить исходную область определения. В этих примерах знаменатель всегда положителен, поэтому посторонние значения $x$ не возникают.

\clearpage
\thispagestyle{fancy}
\begin{center}
{\Large\color{darkaccent}Алгоритм: нахождение ОДЗ}\par
\vspace{5mm}
\end{center}
\begin{center}
\begin{tikzpicture}[
  >={Latex[length=2.4mm]},
  startstop/.style={ellipse,draw=accent,thick,minimum width=4.2cm,minimum height=9mm,align=center,fill=white,font=\normalsize},
  decision/.style={diamond,draw=accent,thick,aspect=2.3,align=center,text width=4.0cm,inner sep=2pt,font=\small},
  action/.style={rectangle,rounded corners=2pt,draw=accent,thick,align=center,text width=4.8cm,minimum height=12mm,inner sep=5pt,font=\normalsize},
  arr/.style={->,thick,color=darkaccent}
]
\node[startstop] (start) at (0,0) {Исходное выражение};
\node[decision] (root) at (0,-1.8) {Есть корни\\ чётной степени?};
\node[action] (rcond) at (-4.0,-3.5) {Для каждого корня: $A(x)\geqslant0$};
\node[decision] (den) at (0,-5.2) {Есть знаменатели?};
\node[action] (dcond) at (-4.0,-6.9) {Для каждого знаменателя: $B(x)\ne0$};
\node[action] (system) at (0,-8.8) {Объедините все условия\\ и решите их};
\node[startstop] (end) at (0,-10.55) {Запишите множество $D$};
\draw[arr] (start)--(root);
\draw[arr] (root.west)-|node[pos=.22,above,font=\small]{да} (rcond.north);
\draw[arr] (rcond.south)|-(den.west);
\draw[arr] (root.east)--++(3.0,0) node[midway,above,font=\small]{нет} |- (den.east);
\draw[arr] (den.west)-|node[pos=.22,above,font=\small]{да} (dcond.north);
\draw[arr] (dcond.south)|-(system.west);
\draw[arr] (den.east)--++(3.0,0) node[midway,above,font=\small]{нет} |- (system.east);
\draw[arr] (system)--(end);
\end{tikzpicture}
\end{center}
\vspace{5mm}
\begin{center}
\color{muted}\small Если ограничений нет, область определения — все действительные числа.\par
Ограничения исходного выражения проверяются совместно.
\end{center}

\clearpage
\sectionhead{9}{Тренировочный блок — Путь А}
Во всех заданиях, где требуется ОДЗ, запишите ответ в виде промежутков или множества чисел. При решении уравнений проверьте найденные корни.

\subhead{Базовый уровень}
\begin{enumerate}
\item Найдите область определения выражения $x^2-9$.
\item Найдите область определения $\sqrt{x^2-16}$.
\item Найдите область определения $\sqrt{25-x^2}$.
\item Найдите область определения $\displaystyle\frac{3}{x^2-4}$.
\end{enumerate}
\subhead{Уровень выше среднего}
\begin{enumerate}[resume]
\item Найдите область определения $\displaystyle\frac{\sqrt{x^2-9}}{x-4}$.
\item Найдите область определения $\displaystyle\frac{1}{\sqrt{x^2-1}-2}$.
\item Решите уравнение $\sqrt{x^2-1}=3-x^2$.
\item Найдите множество значений функции $f(x)=\dfrac{4x}{x^2+4}$.
\end{enumerate}
\subhead{Повышенный уровень}
\begin{enumerate}[resume]
\item Найдите множество значений функции $f(x)=\dfrac{2x+4}{x^2+5}$.
\item Найдите область определения выражения
\[
\frac{\bigl(\sqrt{x^2-9}-2\bigr)\bigl(\sqrt{x^2-9}+1\bigr)}{x^2-11+\sqrt{x^2-9}}.
\]
\end{enumerate}

\clearpage
\sectionhead{10}{Ответы}
\renewcommand{\arraystretch}{1.5}
\begin{tabularx}{\linewidth}{@{}p{10mm}X@{}}
\toprule
\textbf{№} & \textbf{Ответ} \\
\midrule
1 & $\R$. \\
2 & $(-\infty;-4]\cup[4;+\infty)$. \\
3 & $[-5;5]$. \\
4 & $\R\setminus\{-2;2\}$. \\
5 & $\bigl((-\infty;-3]\cup[3;+\infty)\bigr)\setminus\{4\}$. \\
6 & $\bigl((-\infty;-1]\cup[1;+\infty)\bigr)\setminus\{-\sqrt5;\sqrt5\}$. \\
7 & $x=\pm\sqrt2$. \\
8 & $[-1;1]$. \\
9 & $\left[-\frac15;1\right]$. \\
10 & $\bigl((-\infty;-3]\cup[3;+\infty)\bigr)\setminus\{-\sqrt{10};\sqrt{10}\}$. \\
\bottomrule
\end{tabularx}

\vspace{1.0cm}
\subhead{Итоговая памятка}
\renewcommand{\arraystretch}{1.45}
\begin{tabularx}{\linewidth}{@{}p{0.41\linewidth}X@{}}
\toprule
\textbf{Сигнал в задаче} & \textbf{Действие} \\
\midrule
Квадратный корень & Подкоренное выражение $\geqslant0$. \\
Знаменатель & Знаменатель $\ne0$. \\
Корень равен выражению & Изолируйте корень; правая часть $\geqslant0$; квадрат и проверка. \\
Повторяется $x^2$ & Замена $t=x^2$, учитывайте $t\geqslant0$. \\
Модуль не превосходит числа & Переход к двойному неравенству. \\
Множество значений & Положите $y=f(x)$ и найдите все достижимые $y$. \\
Квадратное уравнение относительно $x$ & Проверьте $D_x\geqslant0$ и особые значения параметра $y$. \\
Несколько ограничений & Пересечение допустимых значений. \\
\bottomrule
\end{tabularx}
\end{document}
